\section{Conclusion}

This paper presented KCH-MedRank, a lightweight and reproducible knowledge-constrained local hypergraph learning-to-rank framework for biomedical evidence retrieval and reranking. The method keeps first-stage retrieval practical, constructs a local cross-granularity hypergraph over question, passage, document, entity, MeSH, and relation signals, and uses the resulting medical and structural signals as interpretable learning-to-rank features.

The BioASQ results show stronger top-10 evidence retrieval than original and enhanced Hybrid RRF, with statistically significant improvements across top-10 retrieval metrics against enhanced Hybrid RRF. Against a true MedCPT Cross-Encoder baseline on the same enhanced candidate pool, the strongest supported gain is statistically significant Recall@10 improvement; MRR@10 and nDCG@10 remain positive but non-significant trends. PubMedQA experiments provide a secondary robustness check for evidence coverage, while answer generation remains a secondary task because controlled faithfulness and answer-support metrics do not yet support making it the central contribution.

Overall, the results support a modest claim: knowledge-constrained local hypergraph reranking can improve retrieval of relevant biomedical evidence and provide interpretable evidence signals for evidence-grounded medical QA. The method should be viewed as a retrieval and grounding component, not as a clinically validated answer generation or decision-support system.
